% Vietnamese translation for OI Public Library
% Source-File: IOI中国国家候选队论文/2009/董华星/display/thesis.ppt
\documentclass[11pt]{article}
\usepackage{oipl-vn}

\newcommand{\Oh}{\mathcal{O}}

\title{Tìm hiểu ứng dụng của cây chữ cái trong thi tin học\\{\large Ghi chú theo slide}}
\author{Dong Huaxing\\Trường THPT số 1 Thiệu Hưng, Chiết Giang}
\date{Tháng 12 năm 2008}

\begin{document}
\maketitle

\origtitle{Applications of trie structures in informatics contests}
\sourcepath{IOI-China-Candidate-Team-Papers/2009/Dong-Huaxing/display/thesis-slides.ppt}

\section{Chủ đề}

Slide dùng trie như cấu trúc trung tâm cho nhiều bài xử lý chuỗi và số. Mỗi
đỉnh biểu diễn một tiền tố; đường từ gốc đến đỉnh mô tả khóa đã đọc. Nhờ chia
sẻ tiền tố, nhiều thao tác chuyển từ so sánh chuỗi trực tiếp sang đi trên cây.

\section{Các ứng dụng}

\begin{itemize}
  \item Truy xuất nhanh trên chuỗi: tìm khóa, đếm tiền tố hoặc xác định nhóm
  chuỗi có cùng tiền tố.
  \item Sắp xếp chuỗi: duyệt trie theo thứ tự ký tự để thu được thứ tự từ
  điển.
  \item Giảm chuyển trạng thái vô ích: trong bài có nhiều mẫu, trie gom các
  nhánh giống nhau để tránh kiểm tra lặp.
  \item Tiền tố chung dài nhất: độ sâu của tổ tiên hoặc điểm tách trên trie
  phản ánh độ dài tiền tố chung.
\end{itemize}

\section{Mở rộng}

Trie nhị phân dùng cho số nguyên và các phép xor; trie nén giảm số đỉnh khi
các đường một con quá dài; automaton trên trie hỗ trợ khớp nhiều mẫu. Dù vậy,
trie có thể tốn bộ nhớ nếu bảng chữ cái lớn hoặc dữ liệu ít chia sẻ tiền tố.

\section{Ghi chú}

Theo bài viết đầy đủ, điều quan trọng là nhận ra ``tiền tố'' trong bài toán:
có thể là tiền tố chuỗi, tiền tố bit của số, hoặc tiền tố của một quá trình
quyết định.

\end{document}
